\documentclass[11pt,answers]{exam}
\usepackage[legalpaper,bindingoffset=0in,left=.6in,right=0.8in,top=0.4in,bottom=0.7in,headsep=.5\baselineskip]{geometry}
\usepackage[utf8]{inputenc}
\usepackage{graphicx}
\usepackage{listings}
\usepackage{enumitem}
\usepackage{amsmath}
\usepackage{multirow}
\usepackage{multicol}
\usepackage{caption}
\usepackage{tabularx}
\usepackage{url}
\usepackage{amssymb}
\usepackage{array}
\usepackage{tikz}
\usetikzlibrary{shapes.callouts, positioning, arrows.meta, shapes.geometric, shadows, calc, patterns, 3d, backgrounds, shadings}
\usepackage{adjustbox}
\usepackage{ragged2e}
\usepackage{pgfplots}
\pgfplotsset{compat=1.18} % Added to suppress pgfplots warning

\def\changemargin#1#2{\list{}{\rightmargin#2\leftmargin#1}\item[]}
\let\endchangemargin=\endlist 

\usepackage{xcolor}
\definecolor{mGreen}{rgb}{0,0.6,0}
\definecolor{mGray}{rgb}{0.5,0.5,0.5}
\definecolor{mPurple}{rgb}{0.58,0,0.82}
\definecolor{backgroundColour}{rgb}{0.95,0.95,0.92}
\lstset{
    language=C,
    basicstyle=\ttfamily\small,
    aboveskip={1.0\baselineskip},
    belowskip={1.0\baselineskip},
    columns=fixed,
    extendedchars=true,
    breaklines=true,
    tabsize=4,
    prebreak=\raisebox{0ex}[0ex][0ex]{\ensuremath{\hookleftarrow}},
    frame=lines,
    showtabs=false,
    showspaces=false,
    showstringspaces=false,
    keywordstyle=\color[rgb]{0.627,0.126,0.941},
    commentstyle=\color[rgb]{0.133,0.545,0.133},
    stringstyle=\color[rgb]{01,0,0},
    numbers=left,
    numberstyle=\small,
    stepnumber=1,
    numbersep=10pt,
    captionpos=t,
    escapeinside={\%*}{*)}
}

\pointsdroppedatright

\title{
\vspace{-1.5cm} % Pull the title up closer to the top margin
\begin{changemargin}{0.68cm}{0.68cm} 
\normalsize Class Test 4\hfill {July 2025}\\
\end{changemargin} 
\vspace{.4cm}
\large BANGLADESH UNIVERSITY OF ENGINEERING AND TECHNOLOGY\\
\vspace{.1cm}
Department of Computer Science and Engineering\\
\vspace{.2cm}
Course Code: CSE 219 \hspace{.1cm} Course Title: Signals and Linear Systems\\
\vspace{.2cm}
\textbf{Time: 25 minutes} \hspace{2.2cm} \textbf{Marks: 20}
\begin{changemargin}{0.68cm}{0.68cm} 
\normalsize Student Name: \underline{\hspace{6.5cm}} \hfill Student ID: \underline{\hspace{2.5cm}}\\
\end{changemargin} 
}
\date{} % Prevents today's date from rendering below the title

\begin{document}

\maketitle
\thispagestyle{empty} % Removes the page number from the first page
\vspace{-2.4cm} % Reduces the default space left by maketitle 

\begin{center}
    Answer \textbf{exactly 2 out of Questions 1--4}. \textbf{Question 5 is mandatory.} Answer concisely.
\end{center}

\begin{questions}

    % Question 1
    \question Consider the signal $x[n]$ defined as follows for $N=4$: \hfill [8]
    \[
        x[n] = \begin{cases} 1, & \text{if } n=0 \\ 0, & \text{if } 1 \leq n < N \end{cases}
    \]

    \begin{parts}
        \part Compute the 4-point radix-2 DIT FFT of the signal. Draw a butterfly diagram and show your calculations on the diagram.
        \part If $x[n]$ is viewed as samples of an aperiodic continuous-time signal, what might be a continuous-time signal corresponding to it?
        \part Based on your answer in (b), interpret your result from (a). If the sequence is zero-padded to length $32$, what is its 32-point DFT? Briefly justify.
        \part In a radix-2 DIT FFT, we typically use bit-reversed order for the input $x[n]$. Imagine you forgot to use bit-reversed order and used natural order for the input instead, resulting in a sequence $X'[k]$. Determine the signal $x'[n]$ whose DFT is actually $X'[k]$.
    \end{parts}
    \vspace{0.2cm}

    % Question 2
    \question In many cases, the radix-4 FFT is highly efficient. Suppose you are deriving the radix-4 DIT Cooley-Tukey algorithm. You have already established that for $0 \leq k < N/4$: \hfill [8]
    \[
        X[k] = G_0[k] + W_N^k G_1[k] + W_N^{2k} G_2[k] + W_N^{3k} G_3[k]
    \]
    Where \[
        G_m[k] = \sum_{r=0}^{N/4-1} x[4r+m]\, W_{N/4}^{kr}, \qquad m=0,1,2,3
    \].

    \begin{parts}
        \part Derive the formulas for $X[k+N/4]$, $X[k+2N/4]$, and $X[k+3N/4]$.
        \part Draw the radix-4 butterfly for the case $N=4$. \\
        \textit{Hint: Instead of just $\pm 1$, there will also be multiplications involving $\pm j$.}
    \end{parts}
    \vspace{0.2cm}

    % Question 3
    \question Consider two polynomials: \hfill [8]
    \[
        A(x) = \sum_{n=0}^{N-1} a[n] x^n \quad \text{and} \quad B(x) = \sum_{n=0}^{N-1} b[n] x^n
    \]
    You want to compute their product, $C(x) = A(x)B(x)$.

    \begin{parts}
        \part Provide a simple, naive algorithm to find the polynomial $C(x)$ and analyze its time complexity.
        \part If we express $C(x)$ as the sum $\sum c[n]x^n$, determine the valid range for $n$ and provide a general formula for the coefficients $c[n]$.
        \part Express the coefficient sequence $c[n]$ as a convolution.
        \part Propose a more efficient algorithm to compute $C(x)$ and analyze the time complexity of your proposed algorithm.
    \end{parts}
    \vspace{0.2cm}

    % Question 4
    \question In image processing, we often use a 2-dimensional version of the DFT, which is defined as: \hfill [8]
    \[
        X[k_1,k_2]
        =
        \sum_{n_1=0}^{N_1-1}\sum_{n_2=0}^{N_2-1}
        x[n_1,n_2]\,
        W_{N_1}^{k_1 n_1}
        W_{N_2}^{k_2 n_2}
    \]
    Where $x[n_1, n_2]$ is an $N_1 \times N_2$ matrix and $W_N^k = e^{-j2\pi k/N}$.

    \begin{parts}
        \part What is the time complexity of a naive algorithm used to compute the 2D DFT?
        \part How can you interpret the given formula as two successive 1D DFT operations?
        \part Propose an efficient algorithm (2D FFT) for computing this and analyze its complexity.
        \part What is the primary difference between your 2D FFT algorithm and Bailey's FFT algorithm?
    \end{parts}
    \vspace{0.2cm}

    % Question 5
    \question[4] Thank you for attending the classes and participating in this CT. \hfill [4]

\end{questions}

\end{document}